\documentclass[10pt,a4paper]{article}
\usepackage[utf8]{inputenc}
\usepackage[T1]{fontenc}
\usepackage[margin=1.4cm,top=1.1cm,bottom=1.1cm]{geometry}
\usepackage{amsmath,amssymb}
\usepackage{enumitem}
\usepackage{array}
\usepackage{setspace}
\usepackage{booktabs}
\setstretch{1.0}
\setlength{\parindent}{0pt}
\setlength{\parskip}{0.15em}
\pagestyle{empty}
\newcommand{\repligne}{\par\vspace{0.42cm}\noindent\dotfill}
\newcommand{\zonecalcul}[1]{\par\vspace{0.15cm}\noindent{\small\emph{Espace de calcul :}}\par\vspace{#1}}
\newcommand{\aide}[1]{\fbox{\begin{minipage}{0.96\textwidth}\small #1\end{minipage}}}
\begin{document}
\begin{center}
\textbf{\large Fiche d'exercices --- Séance 9 : apprendre sans étiquettes}\\[0.05cm]
{\small Introduction au ML -- Dr. El Hadji Bassirou TOURÉ -- DMI $\cdot$ FST $\cdot$ UCAD \quad $\bullet$ \quad \textbf{Travail à la main, sans machine}}
\end{center}
\vspace{-0.15cm}\hrule\vspace{0.25cm}
{\small Cette fiche reprend, exercice par exercice, les calculs de la Séance 9 : k-means déroulé à la main (en une puis deux dimensions), l'inertie, le coude et la silhouette, l'effet de l'échelle, puis les fondations de l'ACP --- variance, covariance, vecteurs propres --- et la variance expliquée. Répondre directement sur la fiche.}
\vspace{0.2cm}

\aide{\textbf{Aides numériques (valables pour toute la fiche).} Distance au carré : $\lVert \mathbf x-\mathbf c\rVert^2 = \sum_j (x_j-c_j)^2$. Inertie : $\mathcal I = \sum_k\sum_{i\in C_k}\lVert\mathbf x_i-\mathbf c_k\rVert^2$. Silhouette d'un point : $s = \dfrac{b-a}{\max(a,b)}$ ($a$ : distance moyenne à son cluster, $b$ : au cluster voisin).
Variance $\operatorname{Var}(x)=\frac1n\sum_i(x_i-\bar x)^2$ ; covariance $\operatorname{Cov}(x,y)=\frac1n\sum_i(x_i-\bar x)(y_i-\bar y)$ ; vecteur propre : $A\mathbf v=\lambda\mathbf v$ ; variance expliquée par la composante $i$ : $\lambda_i/\sum_j\lambda_j$.}
\vspace{0.4cm}

\textbf{\large Exercice 1 -- Une étape de k-means en une dimension}
\vspace{0.15cm}

Cinq patients décrits par une seule variable : $x = 2,\ 3,\ 4,\ 10,\ 12$. On vise $k=2$, avec les centres initiaux (mal choisis) $c_1 = 2$ et $c_2 = 3$.

\vspace{0.15cm}
(a) Affecter chaque point au centre le plus proche, puis recalculer les deux centres.
\zonecalcul{1.1cm}
clusters : \makebox[5cm]{\dotfill} \quad $c_1=$ \makebox[1.5cm]{\dotfill} \quad $c_2=$ \makebox[1.5cm]{\dotfill}

\vspace{0.15cm}
(b) Réaffecter avec les nouveaux centres, puis recalculer. Continuer jusqu'à ce que plus rien ne change. Combien d'itérations ?
\zonecalcul{1.4cm}
clusters finaux : \makebox[5cm]{\dotfill} \quad centres finaux : \makebox[3cm]{\dotfill}

\vspace{0.15cm}
(c) Calculer l'inertie avec les centres initiaux ($2$ et $3$), puis l'inertie finale. Que constate-t-on ?
\zonecalcul{1.0cm}
inertie initiale $=$ \makebox[2cm]{\dotfill} \quad inertie finale $=$ \makebox[2cm]{\dotfill}

\vspace{0.45cm}
\textbf{\large Exercice 2 -- k-means à la main, en deux dimensions}
\vspace{0.15cm}

\aide{Six patients (deux mesures exprimées sur une même échelle) :
$A(1,1)$, $B(3,1)$, $C(2,4)$, $D(7,7)$, $E(9,7)$, $F(8,10)$. On vise $k=2$, avec des centres initiaux volontairement médiocres : $\mathbf c_1 = A = (1,1)$ et $\mathbf c_2 = C = (2,4)$.}

\vspace{0.2cm}
(a) \textbf{Affectation 1.} Calculer la distance au carré de chaque point aux deux centres, puis affecter.

\vspace{0.1cm}
{\small \renewcommand{\arraystretch}{1.8}\begin{tabular}{|c|c|c|c|c|c|c|}\hline
point & $A$ & $B$ & $C$ & $D$ & $E$ & $F$\\\hline
$d^2$ à $\mathbf c_1=(1,1)$ & & & & & & \\\hline
$d^2$ à $\mathbf c_2=(2,4)$ & & & & & & \\\hline
cluster & & & & & & \\\hline
\end{tabular}}

\vspace{0.15cm}
Inertie après l'affectation 1 : \makebox[2cm]{\dotfill}

\vspace{0.2cm}
(b) \textbf{Mise à jour 1.} Recalculer les deux centres (moyennes des clusters).
\zonecalcul{0.8cm}
$\mathbf c_1=$ \makebox[2.5cm]{\dotfill} \quad $\mathbf c_2=$ \makebox[2.5cm]{\dotfill}

\vspace{0.15cm}
(c) \textbf{Affectation 2 et mise à jour 2.} Quel point change de cluster ? Donner les nouveaux centres et l'inertie.
\zonecalcul{2.0cm}
point qui change : \makebox[1.5cm]{\dotfill} \quad $\mathbf c_1=$ \makebox[2cm]{\dotfill} \quad $\mathbf c_2=$ \makebox[2cm]{\dotfill} \quad inertie $=$ \makebox[1.8cm]{\dotfill}

\vspace{0.15cm}
(d) Une troisième affectation change-t-elle quelque chose ? Donner l'inertie finale, et dire si k-means garantit toujours la meilleure solution.
\zonecalcul{0.9cm}
inertie finale $=$ \makebox[2cm]{\dotfill}
\repligne

\vspace{0.45cm}
\textbf{\large Exercice 3 -- Combien de clusters ? Coude et silhouette}
\vspace{0.15cm}

(a) Sur un jeu de patients, l'inertie vaut, selon $k$ :

\vspace{0.1cm}
{\small \renewcommand{\arraystretch}{1.3}\begin{tabular}{|c|cccccc|}\hline
$k$ & $1$ & $2$ & $3$ & $4$ & $5$ & $6$\\\hline
inertie & $900$ & $420$ & $180$ & $150$ & $128$ & $110$\\\hline
\end{tabular}}

\vspace{0.1cm}
Calculer le gain d'inertie à chaque ajout de cluster, puis désigner le « coude ».
\zonecalcul{0.9cm}
gains : \makebox[6cm]{\dotfill} \quad coude : $k=$ \makebox[1cm]{\dotfill}

\vspace{0.15cm}
(b) Pourquoi ne choisit-on pas le $k$ qui minimise l'inertie ?
\repligne

\vspace{0.15cm}
(c) Calculer la silhouette de trois points et interpréter chacune.

\vspace{0.1cm}
{\small \renewcommand{\arraystretch}{1.8}\begin{tabular}{|c|c|c|c|c|}\hline
point & $a$ & $b$ & $s$ & interprétation\\\hline
P & $1{,}2$ & $4{,}0$ & & \\\hline
Q & $3{,}0$ & $2{,}5$ & & \\\hline
R & $2{,}0$ & $2{,}2$ & & \\\hline
\end{tabular}}

\vspace{0.45cm}
\textbf{\large Exercice 4 -- k-means et l'échelle des variables}
\vspace{0.15cm}

Trois patients sont décrits par (âge en années, glycémie en g/L) : Awa $(30;\ 1{,}0)$, Binta $(34;\ 2{,}2)$, Coumba $(50;\ 1{,}0)$. Sur le jeu, âge : $\sigma = 10$ ; glycémie : $\sigma = 0{,}5$.

\vspace{0.15cm}
(a) En unités brutes, calculer $d^2(\text{Awa},\text{Binta})$ et $d^2(\text{Awa},\text{Coumba})$. Qui est la plus proche d'Awa ?
\zonecalcul{0.9cm}

(b) Diviser chaque variable par son écart-type (le centrage ne change pas les distances), puis recalculer les deux distances au carré. Qui est maintenant la plus proche ?
\zonecalcul{1.1cm}

(c) Qu'aurait fait k-means sans standardisation sur ce jeu ?
\repligne\repligne

\vspace{0.45cm}
\textbf{\large Exercice 5 -- Variance et covariance}
\vspace{0.15cm}

Deux variables mesurées sur quatre patients : $x = (1,\ 3,\ 5,\ 7)$ et $y = (2,\ 3,\ 7,\ 8)$.

\vspace{0.15cm}
(a) Calculer les moyennes, puis $\operatorname{Var}(x)$, $\operatorname{Var}(y)$ et $\operatorname{Cov}(x,y)$.

\vspace{0.1cm}
{\small \renewcommand{\arraystretch}{1.8}\begin{tabular}{|c|c|c|c|c|c|}\hline
patient & $x_i-\bar x$ & $y_i-\bar y$ & $(x_i-\bar x)^2$ & $(y_i-\bar y)^2$ & $(x_i-\bar x)(y_i-\bar y)$\\\hline
$1$ & & & & & \\\hline
$2$ & & & & & \\\hline
$3$ & & & & & \\\hline
$4$ & & & & & \\\hline
moyenne & --- & --- & & & \\\hline
\end{tabular}}

\vspace{0.15cm}
(b) Écrire la matrice de covariance. Le signe de la covariance était-il prévisible ?
\zonecalcul{0.9cm}
\repligne

\vspace{0.45cm}
\textbf{\large Exercice 6 -- Vecteurs propres d'une matrice de covariance}
\vspace{0.15cm}

Soit la matrice de covariance $C = \begin{pmatrix} 5 & 3\\ 3 & 5\end{pmatrix}$.

\vspace{0.15cm}
(a) Calculer $C\,(1;\ 1)^\top$ et $C\,(1;\ -1)^\top$. En déduire que ce sont des vecteurs propres, et leurs valeurs propres.
\zonecalcul{1.2cm}
$\lambda_1=$ \makebox[1.5cm]{\dotfill} (vecteur \makebox[1.8cm]{\dotfill}) \quad $\lambda_2=$ \makebox[1.5cm]{\dotfill} (vecteur \makebox[1.8cm]{\dotfill})

\vspace{0.15cm}
(b) Quelle est la direction de plus grande variance ? Quelle part de la variance totale porte-t-elle ?
\zonecalcul{0.7cm}
\repligne

\vspace{0.15cm}
(c) Normaliser le premier vecteur propre (longueur $1$). C'est la première composante principale.
\repligne

\vspace{0.45cm}
\textbf{\large Exercice 7 -- La variance expliquée : combien de composantes garder ?}
\vspace{0.15cm}

Une ACP sur six variables standardisées d'un registre de santé donne les valeurs propres :
\[ \lambda = (2{,}4;\ 1{,}2;\ 0{,}9;\ 0{,}7;\ 0{,}5;\ 0{,}3) \qquad (\text{somme} = 6). \]

(a) Compléter la variance expliquée et la variance cumulée.

\vspace{0.1cm}
{\small \renewcommand{\arraystretch}{1.8}\begin{tabular}{|c|c|c|c|c|c|c|}\hline
composante & PC1 & PC2 & PC3 & PC4 & PC5 & PC6\\\hline
part expliquée & & & & & & \\\hline
cumul & & & & & & \\\hline
\end{tabular}}

\vspace{0.15cm}
(b) Combien de composantes garder pour conserver au moins $80\,\%$ de la variance ? Que perd-on en ne gardant que deux composantes pour \emph{visualiser} ?
\repligne\repligne

\vspace{0.15cm}
(c) Pourquoi la somme des valeurs propres vaut-elle exactement $6$ ici ? Pourquoi faut-il standardiser avant une ACP ?
\repligne\repligne

\end{document}
